\capredactadotrue
\section{Descripción General}

\subsection{Descripción básica del negocio}
El proyecto consiste en el desarrollo de una plataforma digital para la gestión y organización de campañas de juegos de rol de mesa (TTRPG, por sus siglas en inglés), en particular de aquellas basadas en sistemas como \textit{Dungeons \& Dragons}. La solución asiste principalmente al \textit{Dungeon Master} (DM) en la planificación y administración de partidas presenciales, centralizando información que habitualmente se encuentra dispersa entre manuales físicos, notas personales y registros individuales de los jugadores.

La plataforma permite crear y gestionar campañas, sesiones y personajes, y registrar eventos, combates, objetos, localizaciones y notas relevantes para el desarrollo narrativo. Los jugadores, por su parte, consultan la información de sus personajes y se incorporan a las campañas mediante invitaciones o códigos de acceso.

A diferencia de los \textit{Virtual Tabletops} (VTT) tradicionales, el proyecto no busca reemplazar la experiencia presencial del juego de mesa, sino complementarla con herramientas de organización, accesibilidad y gestión de la información.

\subsubsection{Tipo de E-Business y modalidad de E-Commerce}
El negocio combina dos modelos de E-Business: una plataforma de software como servicio (SaaS) y un mercado digital (\textit{marketplace}) de contenido. Las modalidades de comercio electrónico son:
\begin{itemize}
  \item \textbf{B2C (empresa a consumidor):} esquema \textit{freemium} orientado al usuario final, con funcionalidades básicas gratuitas y contenido premium personalizable.
  \item \textbf{C2C (consumidor a consumidor):} los creadores comercializan campañas, mapas y otros recursos digitales a otros usuarios dentro de la plataforma, que actúa como intermediaria. Esta modalidad se habilita en la tercera etapa del plan.
\end{itemize}

\subsection{Situación actual del negocio}
El mercado de herramientas digitales para juegos de rol de mesa se encuentra en expansión, impulsado por el crecimiento de la comunidad de jugadores, la popularización del género en medios digitales y el avance de plataformas en línea especializadas. Existen soluciones consolidadas, como Roll20, orientadas principalmente al juego remoto y a la virtualización completa de las partidas.

En este contexto se identifica una oportunidad para herramientas enfocadas específicamente en la gestión de la información y la organización de campañas presenciales, un segmento menos explorado del ecosistema actual.

El proyecto se encuentra en etapa de prototipo validado. Ya se implementaron la creación de campañas y sesiones, la gestión de personajes, la autenticación de usuarios y la participación de jugadores mediante invitaciones. El sistema fue presentado en un entorno académico y validado con potenciales usuarios mediante pruebas y retroalimentación directa. Actualmente se incorporan nuevas funcionalidades: notas públicas y privadas, sincronización de eventos en tiempo real, mapas interactivos y herramientas de asistencia para el seguimiento de las partidas.

\subsection{¿Qué hace único a su negocio?}
La diferencia del proyecto reside en su enfoque como herramienta \emph{complementaria} a la experiencia tradicional de los TTRPG. En lugar de ofrecer un reemplazo digital de la modalidad de juego, prioriza la organización de la información pertinente y agiliza las interacciones más habituales de la mesa. De allí derivan sus principales ventajas competitivas:
\begin{itemize}
  \item un posicionamiento propio, distinto al de los VTT, en el segmento desatendido de la gestión de campañas presenciales;
  \item un núcleo de gestión de la información sin barrera de pago;
  \item una oferta de contenido enriquecida por la comunidad de creadores.
\end{itemize}

\subsection{Factores principales de éxito}
El proyecto busca enriquecer y facilitar la gestión de la información durante la partida, sin las limitaciones ni sustituciones que implica la modalidad completamente digital, considerada la principal barrera de entrada para los jugadores nuevos. Los factores que se consideran determinantes son:
\begin{itemize}
  \item \textbf{Originalidad:} complementa el juego presencial en lugar de virtualizarlo.
  \item \textbf{Precio competitivo:} las funcionalidades centrales de gestión de campañas son gratuitas; el esquema \textit{freemium} reserva para el plan premium el contenido personalizable y los servicios adicionales. Esto otorga mayor competitividad frente a las alternativas del sector.
  \item \textbf{Calidad del servicio:} estabilidad, seguridad y facilidad de uso, condiciones necesarias para la adopción por parte de DMs y jugadores.
  \item \textbf{Experiencia en la actividad:} el equipo conoce el dominio como jugador y cuenta con formación técnica en ingeniería de sistemas.
  \item \textbf{Participación de la comunidad:} el aporte de contenido de los usuarios realza la ambientación del juego y sostiene la cooperación entre jugadores y DMs.
\end{itemize}
Las principales carencias del proyecto son su marca aún desconocida y los recursos limitados de un equipo fundador pequeño. Se abordan mediante la validación continua con usuarios reales, el uso de infraestructura propia de bajo costo, con escalado a la nube a medida que crezca la base de usuarios, y la incorporación gradual de personal.

\subsection{Estrategia: misión, visión y propósito estratégico}
\subsubsection{Misión}
Proveer a jugadores y DMs una plataforma digital unificada que simplifique la organización de la información de las campañas, mejore la planificación de las sesiones y enriquezca la colaboración mediante herramientas accesibles e intuitivas.

\subsubsection{Visión}
Ser una plataforma de referencia en la comunidad de TTRPG, que permita a los jugadores crear experiencias de juego centradas en la creatividad, la inmersión y la cooperación.

\subsubsection{Propósito estratégico}
Permitir que jugadores y DMs dediquen más tiempo al desarrollo de historias, estrategias y experiencias de juego, simplificando las tareas administrativas y centralizando la gestión informática de las campañas.

\subsection{Identificación de la oportunidad de negocio}
En las campañas es frecuente que la información de los jugadores y del DM provenga de fuentes diversas que no se consolidan durante la sesión. Esto perjudica la referencia para planificar sesiones posteriores e incluso el avance de la campaña, con el consiguiente deterioro de la calidad de juego.

Son situaciones habituales la postergación de sesiones por falta de jugadores, la demora en los combates por no contar con información actualizada y la sobrecarga de información en el tablero del DM, que dificulta su participación y atenúa la experiencia de juego.

Los grupos de TTRPG suelen emplear múltiples herramientas aisladas para gestionar campañas, coordinar sesiones y compartir información: hojas en papel, grupos de mensajería o aplicaciones monolíticas. Ello genera fragmentación, duplicación de datos y dificultades de organización. El proyecto busca centralizar estas actividades en una solución integrada, pensada específicamente para las necesidades de jugadores y DMs.

\subsection{Capacidades centrales}
El grupo emprendedor tiene experiencia como jugador de juegos de rol de mesa, con participación activa en campañas y sesiones. Este conocimiento permite comprender las necesidades de los usuarios y las dificultades de organizar campañas, gestionar personajes, coordinar sesiones e intercambiar información entre jugadores y DMs.

En lo técnico, los integrantes cuentan con formación en Ingeniería en Sistemas: análisis, diseño y desarrollo de software, modelado de procesos, diseño de bases de datos, arquitectura de sistemas y metodologías de desarrollo. Estas capacidades permiten diseñar e implementar una solución escalable y alineada con las necesidades identificadas.

La combinación de conocimiento del dominio y capacidades técnicas permite transformar los problemas de la comunidad objetivo en soluciones viables y adaptar la plataforma a nuevas necesidades a lo largo del desarrollo.

\subsection{Propuesta de valor para el cliente}
El proyecto centraliza en un único entorno digital las principales actividades de gestión y participación en campañas de juegos de rol de mesa.

La solución permite a DMs y jugadores administrar campañas, personajes, sesiones e información relevante de manera organizada y accesible. De este modo se reduce la necesidad de recurrir a múltiples herramientas independientes para coordinar partidas, compartir información y seguir los acontecimientos de la campaña.

Además, mejora la experiencia de juego mediante funcionalidades de colaboración, comunicación y acceso rápido a la información, de modo que los usuarios dediquen más tiempo a la narrativa y menos a tareas administrativas y de organización.

\subsection{Valores nucleares de la organización}\label{sec:valores-nucleares}
\begin{itemize}
  \item \textbf{Experiencia de juego:} toda decisión de producto se orienta a mejorar la vivencia de la mesa.
  \item \textbf{Inclusividad:} la plataforma debe ser accesible tanto a jugadores veteranos como a quienes se inician en el género.
  \item \textbf{Colaboración:} el valor de la plataforma crece con la cooperación entre jugadores, DMs y creadores.
  \item \textbf{Inmersión:} las herramientas deben acompañar la narrativa sin interrumpirla.
  \item \textbf{Creatividad:} se fomenta la creación de historias, campañas y contenido propio.
\end{itemize}

\subsection{Enfoque e iniciativas estratégicas}
La organización adopta una estrategia de \emph{enfoque con diferenciación}: se dirige a un segmento específico, el de los grupos que juegan de forma presencial, y ofrece en él una propuesta distinta a la de los VTT. Un núcleo de funcionalidades de gestión gratuito, con valor comparable o mayor al de las alternativas pagas, completa este enfoque sin sacrificar funcionalidades clave y permite mejorar la experiencia con el contenido creado por la comunidad y publicado en la plataforma.

Las iniciativas que impulsan este enfoque son la mejora de la experiencia de juego, la centralización de la información y el crecimiento de la comunidad mediante la colaboración entre usuarios.

\subsection{Áreas claves de resultados}
\begin{itemize}
  \item \textbf{Desarrollo y calidad de la plataforma:} debe garantizar estabilidad, seguridad y facilidad de uso. La experiencia de usuario es determinante para la adopción y permanencia de jugadores y DMs.
  \item \textbf{Crecimiento de la comunidad:} una comunidad activa incrementa el valor de la plataforma mediante la colaboración, el intercambio de contenido y la recomendación entre jugadores.
  \item \textbf{Retención de usuarios:} dado el modelo SaaS \textit{freemium}, es fundamental sostener una alta tasa de uso recurrente que asegure el crecimiento orgánico y la futura conversión a planes premium.
  \item \textbf{Generación y gestión de contenido comunitario:} la participación de creadores es una fuente de diferenciación y crecimiento, pues amplía el catálogo disponible para todos los usuarios.
  \item \textbf{Sostenibilidad económica:} la plataforma debe alcanzar una monetización suficiente, mediante suscripciones premium, comisiones del \textit{marketplace} y servicios adicionales, para sostener el crecimiento operativo.
\end{itemize}

\subsection{Ingreso al sector: estrategias de inserción}
\subsubsection{Estrategia genérica (Porter)}
Según la tipología de estrategias genéricas de \textcite{porter1985}, el proyecto adopta la estrategia de \textbf{enfoque}, combinada con diferenciación. Un liderazgo en costos en todo el mercado no es viable frente a competidores con mayor escala y recursos; en cambio, un segmento acotado (los grupos de juego presencial) permite una oferta especializada que los VTT, orientados a la virtualización completa, no atienden. La gratuidad del núcleo de gestión refuerza este enfoque, pero no lo define.

\subsubsection{Rol competitivo (Kotler)}
Según la clasificación de roles competitivos de \textcite{kotler2016}, la organización ingresa como \textbf{nichera}: una empresa de baja participación que se especializa en un nicho que los líderes atienden de forma marginal. No se prevé confrontar directamente con Roll20, Foundry VTT o D\&D Beyond, sino ocupar un espacio complementario y evitar la competencia frontal.

\subsubsection{Alianzas y \textit{partnerships}}
El posicionamiento se fortalece mediante alianzas con:
\begin{itemize}
  \item \textbf{creadores de contenido}, que alimentan el \textit{marketplace} y amplían el catálogo;
  \item \textbf{comunidades de TTRPG} (por ejemplo, servidores de Discord y grupos locales), que facilitan la captación de los primeros usuarios y la validación de funcionalidades;
  \item \textbf{espacios y eventos del sector} (como tiendas de juegos y convenciones), ámbitos naturales para la presencia de marca en el público presencial.
\end{itemize}
Estas alianzas son una intención de la organización y se detallarán en el plan de marketing.

\subsubsection{Secuencia de inserción}
El ingreso es gradual y se organiza en tres etapas. En la etapa de lanzamiento (primer año) se presenta el producto en Rosario; en la de consolidación (segundo año) se amplía la cobertura a la región y se prepara el \textit{marketplace}; y en la de expansión (tercer año) se alcanza cobertura nacional y se lanza el \textit{marketplace} de contenido de la comunidad. Cada etapa acompaña el crecimiento de la estructura de la organización.
